\chapter{A backend mental model}

A backend is a set of components:
\begin{itemize}[nosep]
  \item \textbf{Edge}: load balancer, gateway, TLS termination
  \item \textbf{Service}: request handlers and business logic
  \item \textbf{Data}: databases and caches
  \item \textbf{Async}: queues and background workers
  \item \textbf{Observability}: logs, metrics, traces
\end{itemize}

\section{Request lifecycle}
A good lifecycle is consistent across endpoints:
\begin{enumerate}[nosep]
  \item authenticate the caller
  \item authorize the action
  \item validate inputs
  \item apply business rules
  \item read/write data
  \item return a stable response shape
\end{enumerate}

\begin{warningbox}{Warning: mixing concerns}
When authorization and validation happen inside business logic, you will duplicate checks.
Put them in the request pipeline, then keep core logic framework-agnostic.
\end{warningbox}

\section{State and side effects}
Most complexity is state and side effects:
\begin{itemize}[nosep]
  \item a database write is a side effect
  \item sending emails is a side effect
  \item charging a payment method is a side effect
\end{itemize}

Backend engineering is controlling these side effects.
